\section*{Working orientation for v\docversion}
\addcontentsline{toc}{section}{Working orientation for v\docversion}
\begin{quote}\small
\noindent\textbf{Reader's map.} The active reading path of the present document is: \emph{orientation and compact front extraction} $\to$ \emph{QAM formal layer} $\to$ \emph{tool and proof-compression layer} $\to$ \emph{open-question / live-burden layer} $\to$ \emph{computational and decoder layer}. The later appendix-style historical material remains available for proof-memory, genealogy, and local rebuild support, but it is not the governing active theorem core.
\end{quote}
This pre-v4 release proceeds directly from the abstract and contents to the nomenclature. The governing current-state claim is the groundwork-complete pre-OP frontier synthesis stated on the title page, in the abstract, and in the current-state audit blocks. Preserved release-roadmap notes, checkpoint transitions, and earlier front-matter control blocks remain part of the document as inherited baseline material for theorem-memory, genealogy, and reconstruction support; they are no longer to be read as the overriding current-working claim of v\docversion.

\paragraph{r06 bibliography hardening.}
A full-document literature scan was performed before this public split-bundle release. The scan found that the mathematical body repeatedly uses Teichmuller and quasiconformal language, Beltrami-style deformation intuition, Riemann--GUE/RMT spectral heuristics, Hopf/fibre-bundle topology, Dixon/octonionic and division-algebraic language, Dirac/Higgs/Yang--Mills and CKM/PMNS landmarks, Julia-boundary diagnostics, ZF/QAM coding language, and probe/cryptographic/computational analogies. Version~\docversion\ therefore adds a broader terminal bibliography and uses this paragraph as the compact external anchor surface: Teichmuller/quasiconformal background is anchored in \cite{Ahlfors2006,GardinerLakic2000,Hubbard2006,Sullivan1985}; Riemann-zeta, spectral, and GUE/RMT background in \cite{Riemann1859,Edwards1974,Titchmarsh1986,Conrey2003,Montgomery1973,Odlyzko1987,Mehta2004,KatzSarnak1999}; Hopf, fibre-bundle, and exceptional-geometry background in \cite{Hopf1931,Steenrod1951,Hatcher2002,Bryant1987,EguchiGilkeyHanson1980}; division-algebraic and octonionic background in \cite{Baez2002,Dixon1994,Furey2018}; QFT/gauge and mixing landmarks in \cite{Dirac1928,YangMills1954,EnglertBrout1964,Higgs1964,Cabibbo1963,KobayashiMaskawa1973,MakiNakagawaSakata1962}; complex-dynamical Julia background in \cite{Milnor2006}; ZF/set-theoretic foundations in \cite{Zermelo1908,Fraenkel1928,Jech2003,Kunen1980}; and the computational/probe/cryptographic side in \cite{Turing1936,Kasiski1863,RivestShamirAdleman1978,DiffieHellman1976,Norris1997,Glasserman2004,Bollerslev1986,Knuth1997,Rojas1998,Copeland2006}.

\paragraph{r08 MML/MMA compression map.}
The r07 local language interface is now used as a conservative compression and audit map for inherited HC/OP text. The purpose is not to delete historical proof material, but to mark repeated structural paragraphs as instances of typed operations such as \(\mathrm{CAPTURE}_{\mathrm{HC}}\), \(\mathrm{HANDOFF}_{\mathrm{HC}\to\mathrm{OP}}\), \(\mathrm{NORMALIZE}_{\mathrm{OP}}^{\mathrm{HC}}\), \(\mathrm{ROUTE}_{\mathrm{OP}}^{\mathrm{HC}}\), \(\mathrm{COMPARE}_{\mathrm{HC}}^{1/4\mid 5B}\), \(\mathrm{RESOLVE}\), \(\mathrm{RETURN}_{\mathrm{OP}}\), \(\mathrm{SCAN}\), \(\mathrm{RANK}\), and \(\mathrm{CERTIFY?}\). The compression map is therefore a reading and refactoring guide: it can shorten later exposition and make inherited blocks easier to audit, but it cannot manufacture certificates or remove the historical witness.

\paragraph{r10 first tool-contract application.}
The r09 tool library is now applied back to the active HC/OP proof spine. The purpose is deliberately modest: repeated paragraphs in the r01--r04 chain may be referenced by the contracts \(\mathfrak T_{\mathrm{lock}}\), \(\mathfrak T_{\mathrm{handoff}}\), \(\mathfrak T_{\mathrm{norm}}\), \(\mathfrak T_{\mathrm{route}}\), \(\mathfrak T_{\mathrm{cmp}}\), \(\mathfrak T_{\mathrm{resolve}}\), \(\mathfrak T_{\mathrm{cert}}\), and \(\mathfrak T_{\mathrm{frontier}}\) once their local hypotheses and guards have been stated. This is a citation and audit layer, not a shortening pass that deletes proof content.


\paragraph{r11 guarded reduction protocol.}
The present revision adds the first reduction discipline for using those tool contracts in future text. A paragraph is reducible only when it repeats an already stated safety, route, frontier, diagnostic, or no-false-certificate clause. It is not reducible when it introduces a new object, a first hypothesis, a local guard, a certificate test, a proof restriction, a theorematic status change, or historical witness material. Thus the r11 layer is the first practical simplifier produced by MML/MMA, but it remains conservative: it turns repeated prose into guarded citations and never turns an execution trace into an OP closure certificate.


\paragraph{r12 guarded-reduction application pass.}
The present revision applies the r11 discipline for the first time without deleting text. It installs a finite ledger of guarded citation sites along the active HC/OP proof spine: lock/capture reminders, handoff and normalization reminders, route-frontier reminders, paired-comparison status clauses, probe diagnostic clauses, and archive-to-tool extraction reminders. Each entry keeps its first occurrence and local guard explicit, records the relevant tool contract, and permits only later repeated tails to be cited in guarded normal form. Thus r12 is a practical audit layer: it tells the reader exactly where the new compression technology helps the existing text, while preserving theorematic status and the historical witness.


\paragraph{r13 guarded-reduction pilot normal form.}
The present revision turns the r12 ledger into a first usable rewriting pattern. It does not shorten the historical witness and does not replace first definitions. It only says that, after a first occurrence has stated its hypotheses and local guard, a later repeated HC/OP tail may be written in the guarded citation form
\[
  \mathrm{GRCite}_{\mathrm{HCOP}}(\ell;G,W,S),
\]
where \(\ell\) names the ledger entry, \(G\) points to the local guard, \(W\) points to the preserved witness or first occurrence, and \(S\) records the non-certificate status. This is the first concrete simplification normal form produced by the MML/MMA layer: it compresses repeated prose, not proof content.

\paragraph{r14 tail-normalized proof skeleton.}
The present revision organizes the r13 guarded-citation mechanism into a first tail-normalized proof skeleton for the active HC/OP spine. Its purpose is to show where the first explicit mathematical source block remains, where later repeated endings may be cited in guarded form, and how the dependency order of lock/capture, handoff, normalization, routing, paired comparison, and resolution is preserved. Thus r14 does not shorten the document by deletion. It gives future HC/OP text a compact reference skeleton in which repeated tails can be replaced by guarded citations only after their source block, local guard, witness, and theorematic status have already appeared.
\paragraph{r15 skeleton-to-workstep dispatch.}
The present revision turns the r14 tail skeleton into a first controlled dispatch object. Its purpose is not to claim an OP closure, but to say which finite next work packets may be launched from the ordered skeleton nodes. The admissible dispatch order remains lock/capture, handoff, normalization, routing, comparison, and resolution; each dispatched packet must keep its local guard, witness, route tag, and non-certificate status unless a separate explicit certificate test is later supplied.

\paragraph{r16 first dispatched lock/capture stabilization.}
The present revision executes the first dispatched work packet from the r15 menu. It acts only on the LOCK/CAPTURE node, records a finite lock-stability defect register, and returns one conservative status: handoff-ready stable source, bounded local repair source, or smaller frontier source. This is the first point at which the skeleton-to-workstep dispatch becomes an actual controlled mathematical work step; it still does not generate an OP closure certificate.


\paragraph{r07 QAM/MML/MMA local interface registry.}
Version~\docversion\ records the narrow local interface stamps for the machine-language layer in prose and uses symbolic local interface tags in formulas for the current HC-routed OP-entry chain. The historical computational chapter already contains the broader MML--MMA--MCP reserve architecture. The present r07 step does not activate that reserve as an independent proof engine. It only extracts the finite sublanguage needed to read the r01--r04 chain as a typed program:
\[
  \mathrm{QAM}_{\mathrm{HCOP}}^{\mathrm{loc}}
  \longrightarrow
  \mathrm{MML}_{\mathrm{HCOP}}^{\mathrm{loc}}
  \longrightarrow
  \mathrm{MMA}_{\mathrm{HCOP}}^{\mathrm{loc}}.
\]
Here \(\mathrm{QAM}_{\mathrm{HCOP}}^{\mathrm{loc}}\) is the local bridge tag over the inherited QAM-over-ZF object language, with interface stamp QAM-HC/OP 0.1 recorded as metadata; \(\mathrm{MML}_{\mathrm{HCOP}}^{\mathrm{loc}}\) is the finite instruction interface for capture, handoff, normalization, routing, comparison, and return, with interface stamp MML-HC/OP 0.2 recorded as metadata; and \(\mathrm{MMA}_{\mathrm{HCOP}}^{\mathrm{loc}}\) is the route-faithful execution semantics for those instructions on finite HC/OP states, with interface stamp MMA-HC/OP 0.2 recorded as metadata.


The following inherited front block should therefore be read as the carried-forward active baseline from the audited v3.30.0 line, with the following priorities and outcomes:
\begin{enumerate}[leftmargin=2em]
  \item \textbf{Preserve the stabilized closure/interface baseline.} The language/closure, branch/shell/routed, tensor, operator-family, translation, OP~3 shell-certificate, bundled interface compatibility, and terminal proof-use surface already exported in earlier public stages are treated as fixed baseline rather than reopened here.
  \item \textbf{Keep the partial-closure/no-second-carrier layer fixed.} The terminal-interface no-second-carrier theorem, the first partial OP~3/OP~5 closure theorem, and the restricted lift-back principle are treated as available structural infrastructure.
  \item \textbf{Keep the first retained-carrier optical approximation and machine-shadow baseline fixed.} Reduced prism readout, residual diagnostics, carrier-side fit/calibration/validation, inversion, refinement, reduced stability, and the bidirectional observer-side split \(\eta_+=\eta_{\mathrm{drive}}+\eta_{\mathrm{carrier}}\), \(\eta_-=\eta_{\mathrm{drive}}-\eta_{\mathrm{carrier}}\) are treated as the stabilized computational surface rather than rebuilt from scratch.
  \item \textbf{Inherit the completed tool/QAM/decoder layer as fixed infrastructure.} Channel-sorted total, imbalance, and bridge packets, finite QAM submenus, primary tags, registry assignments, tag-first reopening routes, the Monte-Carlo residual-signature routing layer, the action table, deterministic decoder syntax, transition function, observer/event normal form, evaluation map, and compact output tokens are treated as the completed route-control package inherited from v3.28.0.
  \item \textbf{Keep the prior post-closure burden reading visible as inherited baseline.} The v3.29.0--v3.30.0 burden picture is preserved here as inherited local baseline material from which the new pre-OP frontier synthesis is now read, not as an independent current public-claim layer that overrides the new groundwork status stated in the opening pages.
  \item \textbf{Keep the DOI logic normalized.} The concept DOI remains the stable title-page pointer of the Companion line. The governing latest public-release DOI in the chronology is now \href{https://doi.org/10.5281/zenodo.19750674}{10.5281/zenodo.19750674} for the public v3.35.0 release, while \href{https://doi.org/10.5281/zenodo.19746821}{10.5281/zenodo.19746821} remains the prior public release DOI for v3.34.0.
\end{enumerate}

\paragraph{Reading rule for the inherited active-front block.}
The inherited v3.30.0 front layer is kept intact as carried-forward baseline source material. It records the then-active retained-carrier optical and live-burden reading, including OP~1 and OP~5 Step~(B) closure language on the inherited admissible corridors, and should not be read as the governing present-release claim of v\docversion. In the present release final, those statements are preserved for theorem-memory and genealogy; the governing current-state claim, however, remains the frontier-certified pre-entry synthesis developed in the new post-v3.31.0 line. Accordingly, references below to live OP burden, route sharpening, kernel closure, or routed-balance obstruction inside the inherited front block should be read as preserved baseline source material rather than as the final current release claim of this release.

\subsection*{Compact public-shape branch/shell/routed extraction}
\begin{quote}\small
\noindent\textbf{Status.} \emph{Active front-facing compression block.}\par
\noindent\textbf{Block function.} This block gives the smallest front-facing statement of the public branch/shell/routed QAM package, so that a reader can see what is already stabilized before entering the fuller internal unfolding.
\end{quote}
The inherited public-shape QAM layer may be stated in compressed form before the full internal unfolding is read. In that form, one starts with the already public coded source/readout, branch/shell/routed, tensor, operator-family, and translation layers from v3.15.0--v3.17.0, then adds shell certificates, shell-defect witnesses, kernel-persistence statements, and shell-stable readout factorization. The guiding claim is that these added roles do not create a new closure criterion; they only make the OP~3 shell layer readable as one explicit coded closure mechanism.

\begin{definition}[Compact public-shape QAM package]
The \emph{compact public-shape QAM package} for the inherited public-shape v3.16.0 release layer consists of the following release-level data:
\begin{enumerate}[leftmargin=2em]
  \item tagged coded objects and minimal predicates for source states, layers, operators, readouts, branches, shell/refinement objects, admissibility witnesses, and tensor-coded composites;
  \item admissible coded transport together with coded closure-class preservation and closure-compatible coded readout;
  \item closure-admissible branch codes;
  \item shell-admissible refinement codes over closure-admissible branches;
  \item normalized admissible routed support over the admissible branch-shell support class.
\end{enumerate}
At this public-shape stage, tensor admissibility, operator-family/lambda-like maps, and support-selection or degeneracy refinements are not part of the compact release claim, even though they remain internally available below.
\end{definition}

\begin{theorem}[Compact public-shape preservation chain]\label{thm:compact-public-shape-preservation}
Let $s_{\mathrm{in}}$ be a coded source state, let $b$ be a closure-admissible branch code, let $\sigma$ be a shell-admissible refinement code on $b$, let $u$ be a normalized admissible routed support code over the admissible branch-shell support class, and let $r$ be a closure-compatible coded readout. Then the v3.16.0 public-shape layer preserves the same closure-relevant readout content along the full compressed chain
\[
s_{\mathrm{in}}
\longrightarrow
(b,\sigma)
\longrightarrow
u
\longrightarrow
r,
\]
provided the coded return realizations are taken inside the admissible branch-shell support sector.
\end{theorem}

\begin{proof}
At OP level, the same compression is built in stages: invariant source-return closure for the OP~1--OP~5 pair is fixed by Theorem~\ref{thm:op15-closure}; branchwise preservation is then isolated by Proposition~\ref{prop:op-branchwise-refinement}; shell-sensitive preservation at the OP~3 interface is isolated by Theorem~\ref{thm:op3-shell}; and normalized routed preservation is isolated by Proposition~\ref{prop:op-admissible-routing}. The QAM translation theorem (Theorem~\ref{thm:qam-translation}) carries exactly that OP chain into coded form. Later in the QAM section this becomes the admissible branch-shell theorem (Theorem~\ref{thm:qam-branch-shell}), the routed admissible support theorem (Theorem~\ref{thm:qam-routed-support}), and the visible/structural separation principle (Corollary~\ref{cor:qam-visible-structural}). The present statement records that entire OP-to-QAM compression as one front-facing release theorem.
\end{proof}

\begin{remark}[Public-shape reading guide]\label{rem:public-shape-reading-guide}
For the purposes of the current release, the later QAM section should be read with the following dictionary in mind:
\begin{enumerate}[leftmargin=2em]
  \item ``coded source state + closure-compatible coded readout'' is the coded front-face of the OP~1/OP~5 source-return pair;
  \item ``closure-admissible branch code'' is the coded branchwise refinement role;
  \item ``shell-admissible refinement code'' is the coded OP~3 shell role;
  \item ``normalized admissible routed support'' is the coded routed aggregation role over the admissible branch-shell support sector;
  \item ``closure-relevant readout content'' always means readout agreement after passage through the associated closure-factor map.
\end{enumerate}
This dictionary is the intended reading bridge between the compact front theorem and the later theorem-level unpacking in Theorem~\ref{thm:qam-branch-shell}, Theorem~\ref{thm:qam-routed-support}, Theorem~\ref{thm:qam-translation}, and Corollary~\ref{cor:qam-visible-structural}.
\end{remark}

\begin{corollary}[Release-level visible/structural rule]\label{cor:release-visible-structural-rule}
Within the compact public-shape QAM package, visible coded mismatch in packet, phase-lock, or $\varepsilon$-screening data is not by itself a closure failure. At the inherited public-shape level, genuine closure failure begins only when branch admissibility or shell admissibility is lost.
\end{corollary}

\begin{remark}[Structural provenance of the public-shape extraction]
The v3.16.0 front block should be read as the coded shadow of one already stabilized OP chain and not as a freestanding new criterion. In that reading, coded source state and coded readout come from the closed OP~1--OP~5 source-return pair of Theorem~\ref{thm:op15-closure}; closure-admissible branch codes come from the branchwise refinement layer of Proposition~\ref{prop:op-branchwise-refinement}; shell-admissible refinement codes come from the OP~3 shell theorem of Theorem~\ref{thm:op3-shell}; normalized admissible routed support comes from the routed aggregation layer of Proposition~\ref{prop:op-admissible-routing}; and the release-level visible/structural rule is the QAM-over-ZF front-face of the OP screening logic, where packet/phase/epsilon data remain preselection data unless they witness loss of admissibility. The public-shape extraction therefore does not add a second closure notion. It republishes the already existing source-return closure architecture in the smaller coded alphabet needed for the current QAM-over-ZF line.
\end{remark}

\begin{remark}
This compact block is intentionally stronger in readability than in scope. It exposes the third public-shape QAM step as one coherent public claim, but it does not yet export the full internal support-selection or wider compression machinery. In other words, v3.17.0 was prepared as the controlled public extraction of the tensor/operator layer, not as the publication of the entire internal QAM architecture.
\end{remark}

\begin{remark}[Deferred kernel-first outlook]
The front theorem should also be read as the visible place where any later zero-kernel migration would have to pay off. The deferred kernel notes below do not introduce a second active tuple language, but they already identify the only later proof-compression route that would matter on the present public surface: first kernel-first transport reduction (Proposition~\ref{prop:qam-kernel-first-transport}), then local branch-shell kernel reduction (Proposition~\ref{prop:qam-kernel-first-branch-shell}), and only afterwards routed family compression (Proposition~\ref{prop:qam-kernel-first-routed}). Read together with the roadmap-fit remark, the later checkpoint summary, the activation criterion of Theorem~\ref{thm:qam-zero-kernel-activation-criterion}, and the later remarks on the staged activation roadmap, the criterion-to-roadmap reading rule, and the compressed synthesis of the deferred zero-kernel path, this deferred path should therefore be understood as a later proof-theoretic refinement of the same front chain rather than as a competing current-release line.
\end{remark}

\begin{remark}[QAM-over-ZF entry layer]
A next formal development line of the Companion is the introduction of QAM not as a replacement for Zermelo--Fraenkel set theory, but as a structure-first coded language over ZF. The point of this layer is to give the current OP/branch/shell architecture a cleaner formal interface for states, layers, operators, readouts, admissibility, closure preservation, and tensor-coded composition while remaining, in principle, back-translatable into set-theoretic form.

At the present integrated stage, the admissible OP fragment remains the first translation target of this QAM line: the OP~1/OP~5 source-return closure architecture, admissible branch transport, OP~3 shell-admissible refinement, admissible routed support, tensor composition over admissible coded objects, and the visible packet/phase-lock/epsilon screening layer. Relative to the public v3.16.0 release, the next public-shape export in that historical sequence was aimed at the tensor and operator/lambda layer rather than at a further broadening of the branch/shell/routed block.

The reserved minor-release path remains the following. First, v3.15.0 fixed the minimal QAM language layer over ZF, including tagged set-codes for states, layers, operators, readouts, branches, shell/refinement objects, admissibility witnesses, and tensor-coded composites. Second, v3.16.0 extracted the branch/shell/routed translation into public-shape form, including closure-admissible branch codes, shell-admissible refinement codes, admissible routed support, and the first coded branch-shell transport results. Third, v3.17.0 realized the tensor/operator public line; fourth, v3.18.0 exported the first public OP~3/QAM closure package; fifth, v3.19.0 hardened the witness-grade reopening discipline of that bundled channel; and the present v\docversion\ release line now pilots a closure-liquid-crystal / holographic Julia-boundary re-reading above that compressed surface while support-selection refinements remain below as later-stage technical support.

The present integrated continuation therefore exposes the compressed public-shape QAM core together with its branch/shell/routed baseline, its tensor/operator extension, its public OP~3/QAM closure package, its bundled witness-grade compression surface, and the released v3.20.0 crystal--shadow architecture: tagged coded objects and minimal predicates, admissible transport with coded closure-class preservation, closure-compatible coded readout, closure-admissible branch codes, shell-admissible refinement codes, admissible routed support, tensor admissibility, tensor closure, first operator-family or lambda-like translation statements, terminal proof-use compression, addressed crystal packets/corridors/classes, theorem-level shadows, and the terminal crystal--shadow interface. The broader internal QAM package remains available below as technical support and does not alter the mathematical claims already frozen in the public v3.20.0 release.

\textbf{QAM-coded object preview.} The minimal intended object types are states, layers, operators, readouts, branches, shell/refinement objects, admissibility/closure witnesses, and tensor-coded composites. A set $x$ is intended to count as a QAM-coded object exactly when it has one of the tagged forms
\[
\begin{aligned}
&\langle 0,a\rangle,\quad
\langle 1,\ell,D\rangle,\quad
\langle 2,\ell_1,\ell_2,G\rangle,\quad
\langle 3,\ell,W,R\rangle,\\
&\langle 4,b,B\rangle,\quad
\langle 5,r,s,T\rangle,\quad
\langle 6,u\rangle,\quad
\langle 7,p,q\rangle.
\end{aligned}
\]
If $p$ and $q$ are QAM-coded objects, the tensor-coded composite is written as
\[
p \otimes q := \langle 7,p,q\rangle.
\]
At the entry stage, these tags and tensor codes are only coding devices that distinguish structural roles inside ZF and reserve tensor syntax for later internal hardening.
\end{remark}

\begin{remark}[Dependency guide for the internal QAM package]
The internal QAM-over-ZF line should be read in the following order: coded objects and predicates first, admissible transport and closure-class schemata second, branch/shell/routed admissibility third, tensor admissibility and tensor closure rules fourth, faithful translation of the admissible OP fragment fifth, operator-family/lambda-like structural maps sixth, and support-selection/degeneracy handling only as a seventh-layer refinement on top of the already admissible fragment.
\end{remark}

\subsection*{How to read the internal QAM package}
The internal QAM package should be read in two layers. At the compact reader level, one should keep only the following chain in view: coded objects and predicates, admissible transport and closure, branch/shell/routed admissibility, tensor closure rules, faithful translation of the admissible OP fragment, and only afterwards operator/lambda and support-selection refinements. At the technical level, the detailed sections below provide the full internal unfolding of these same layers in theorematic form.

\begin{remark}[QAM normalization guide]
The current internal package is not meant to be enlarged indiscriminately. Its next intended use is normalization and later public extraction. In particular, later work should prefer compression, dependency control, and clean appendix-style grouping over uncontrolled addition of parallel QAM subsystems.
\end{remark}

\subsection*{Roadmap and extraction criteria for the v3.17.0 public-shape preparation}
The internal QAM package is now far enough along that one may formulate a controlled extraction path toward the next public QAM-oriented line. The intended purpose of this extraction is not to publish the entire internal QAM package at once, but to expose its stable and readable core in stages.

The current preferred extraction sequence is the following.
\begin{enumerate}[leftmargin=2em]
  \item \textbf{First public QAM line (realized in v3.15.0): language and closure interface.} Publish a compact QAM-over-ZF entry layer together with the minimal predicate package, admissible transport, coded closure-class preservation, closure-compatible readout, and a first translation statement for the admissible OP fragment.
  \item \textbf{Second public QAM line (realized in v3.16.0):} \textbf{branch/shell/routed translation.} Extend the public QAM line by adding closure-admissible branch codes, shell-admissible refinement codes, admissible routed support, and the first coded branch-shell transport results.
  \item \textbf{Third public QAM line (realized in v3.17.0): tensor and operator/lambda extension.} Add tensor admissibility and tensor closure rules together with the first operator-family or lambda-like structural maps and explicit translation or round-trip statements, while keeping selection/support handling internal for now.
  \item \textbf{Fourth public QAM line (realized in v3.18.0): first public OP~3/QAM closure package.} Expose shell certificates, shell-defect localization, kernel persistence under the bridge hypotheses, shell-stable readout factorization, and the bundled OP~3$\to$OP~5/global interface in compact public shape.
  \item \textbf{Fifth public QAM line (realized in v3.19.0): witness-grade hardening and local reopening discipline.} Add witness-grade proof-use control for the bundled channel together with a first narrow licensed kernel-first rewrite window, a first corridor-level activation lemma, a first corridor-to-bundle propagation rule, a first bundle-level reopening compression principle, a first profile-to-run compression principle, a first word-determined run-class step, a first run-class concatenation principle, a first finite run-class normal-form principle, a first normal-form-determined bundled proof-use class step, a first proof-use-class composition principle, a first associative proof-use-class normal-form principle, a first contraction-stable class-normal-form principle, and a terminal proof-use surface.
  \item \textbf{Sixth public QAM line (realized in v3.20.0): closure-liquid-crystal, addressed role layer, theorem-level shadows, and terminal crystal--shadow interface.} Re-read OP~3 as local cell compatibility, OP~5 as defect transport on the same cellular carrier, visible sectors as liquid-crystal domains of one common bulk order, Julia-type diagnostics as holographic boundary readout, stabilize a first base36/Q-word role addressing with reserved 0-band transport support, and connect the resulting terminal crystal-readout surface to a bundled theorem-level shadow package and a terminal crystal--shadow interface.
\end{enumerate}

At the present internal stage, the extraction criteria for the sixth public QAM line should remain strict. A QAM block should be promoted into a public release only if it satisfies all of the following conditions:
\begin{enumerate}[leftmargin=2em]
  \item \textbf{compression:} the public block can be read in a compact main-text form without forcing the full internal package into the public line;
  \item \textbf{dependency clarity:} the block has a clean dependency guide and does not rely on hidden internal lemmas;
  \item \textbf{ZF-groundedness:} the coded language remains explicitly readable as a definitional overlay over ZF rather than as an unexplained new ontology;
  \item \textbf{translation contact:} the public block makes clear what part of the current OP/branch/shell architecture it actually translates;
  \item \textbf{non-overreach:} the public block does not yet claim more than the present coded layer genuinely supports.
\end{enumerate}

\begin{remark}
The current internal package already satisfies much of this extraction logic, but it is still broader than what should be exported immediately. The next internal passes should therefore favor tensor/operator compression, translation clarity, and clean staging over further rapid enlargement.
\end{remark}

\begin{remark}[How the present release-pass checkpoint fits the public-shape roadmap]\label{rem:qam-roadmap-checkpoint-fit}
The deferred kernel-first checkpoint developed later in this file should be read as an internal support block for the same v3.17.0 public-shape preparation path, not as a second release claim. Its current function is to prepare a later proof-compression interface for the compact tensor/operator chain by isolating the future kernel vocabulary, bridge data, staged stress tests, and activation criteria needed for a controlled migration. In particular, the checkpoint helps explain why the third public QAM line should remain narrow at visible level while still allowing later internal compression work to accumulate behind it.
\end{remark}

\subsection*{Compressed public-shape QAM branch-shell core}
The second public QAM line should expose only the smallest next coded layer needed to carry the current admissible OP fragment and its branch/shell/routed refinement into a ZF-grounded formal interface. At compressed level, this public-shape branch-shell core should be read through six ingredients:
\begin{enumerate}[leftmargin=2em]
  \item tagged QAM-coded objects and the minimal predicate layer;
  \item admissible transport together with coded closure-class preservation;
  \item closure-compatible coded readout;
  \item closure-admissible branch codes;
  \item shell-admissible refinement codes;
  \item admissible routed support together with first coded branch-shell transport statements.
\end{enumerate}

\begin{theorem}[Compressed public-shape QAM branch-shell theorem]
Assume that coded QAM objects are read as a definitional overlay over ZF, that admissible coded transport preserves coded closure class, and that the relevant coded readout is closure-compatible. Assume further that closure-admissible branch codes and shell-admissible refinement codes preserve coded closure class, that normalized admissible routed supports use only such branch-shell pairs, and that the admissible OP fragment is translated faithfully into the coded QAM layer. Then the translated source/return closure content of the admissible OP fragment, together with its admissible branchwise, shellwise, and routed realizations, is preserved at coded level, and closure-relevant readout agreement is retained after passage through the closure-factor map.
\end{theorem}


\begin{remark}
This compressed theorem is the intended mathematical core of the second public QAM line. It remains deliberately weaker than the full internal QAM package: tensor rules, operator/lambda maps, and support-selection layers remain available internally, but they need not all be exported at once in order to make the next public QAM step mathematically meaningful.
\end{remark}

\subsection*{Structural systems view: QAM as a virtual machine with a ZFS-style theorem store}
\begin{quote}\small
\noindent\textbf{Status.} \emph{Active explanatory overlay, not a second mathematical criterion.}\par
\noindent\textbf{Block function.} This block explains why the QAM objects, kernel states, snapshots, and safe side branches belong to one architecture, without changing the theorematic content.
\end{quote}
The present QAM export should also be read through a systems-level analogy. The visible coded tuple language plays the role of a readable virtual-machine interface over the ZF-grounded base, while the later exposed payload kernels and kernel normal forms behave like the machine-state identities that remain stable beneath several visible tuple surfaces. In that sense, the public QAM line is the user-facing instruction and state surface, whereas the deeper kernel-first line is the later OS-level state-management discipline. This reading is not introduced as a replacement for the mathematics, but as a compact way to explain why branch/shell/routed transport, kernel normal forms, search keys, snapshots, and safe side branches fit together as one contiguous architecture. It also aligns with the earlier Zuse Z3 comparison already used in the project: the current point is to interpret the coded QAM layer not merely as notation, but as a structured machine view with auditable state persistence.

\begin{remark}[VM/OS reading of the current QAM line]\label{rem:qam-vm-os-reading}
For the present document, the systems analogy should be read in five layers.
\begin{enumerate}[leftmargin=2em]
  \item \textbf{Visible QAM views.} The active tuple surfaces such as $\langle 0,a\rangle$, $\langle 0,a,b\rangle$, branch codes, shell/refinement codes, and routed return codes form the readable public interface of the machine.
  \item \textbf{Kernel / VM state.} The exposed payload kernels, their normal forms, and the resulting class-level closure/readout data form the later machine-state identity beneath that visible interface.
  \item \textbf{OS-level orchestration.} Admissible transport, local branch/shell return, and routed family aggregation act like the current system-level transition and scheduling layer that moves visible states while preserving closure-relevant content.
  \item \textbf{ZFS-style theorem store.} Search keys, snapshots, clones, audits, and later proof-compression branches belong to an internal persistence layer that manages state identity, safe side branches, and deduplication of repeated kernel classes.
  \item \textbf{Pool/group structure.} Sphere objects may later be treated as logical devices, while trioles or routed-support families may be grouped as internal device families. At the present stage, mirror-style reuse of equal kernel classes is the clearest fit; parity-like reconstruction remains exploratory and should be treated only as a later structural option.
\end{enumerate}
\end{remark}

\begin{figure}[H]
\centering
\resizebox{0.94\linewidth}{!}{%
\begin{tikzpicture}[
  font=\small,
  mainbox/.style={draw, rounded corners=3pt, align=center, minimum width=8.0cm, minimum height=1.05cm, inner sep=6pt},
  sidebox/.style={draw, rounded corners=3pt, align=center, minimum width=3.9cm, minimum height=1.0cm, inner sep=5pt},
  arr/.style={-{Latex[length=2.6mm]}, line width=0.5pt}
]

\node[mainbox, fill=blue!6]   (visible)  at (0,0)   {\textbf{Visible QAM views}\\Tuple surfaces, branch/shell codes, routed returns};
\node[sidebox, fill=yellow!10] (public)   at (7.0,0)   {\textbf{Public surface}\\Readable theorem interface};
\node[mainbox, fill=green!6]  (kernel)   at (0,-1.9) {\textbf{Kernel / VM state}\\Payload kernels, normal forms, closure/readout identity};
\node[mainbox, fill=orange!8] (os)       at (0,-3.8) {\textbf{OS-level orchestration}\\Admissible transport, branch/shell return, routed aggregation};
\node[mainbox, fill=purple!7] (store)    at (0,-5.7) {\textbf{ZFS-style theorem store}\\Search keys, snapshots, clones, audits, safe side branches};
\node[sidebox, fill=yellow!10] (internal) at (7.0,-5.7) {\textbf{Internal persistence}\\State reuse and branch hygiene};
\node[mainbox, fill=gray!12]  (pool)     at (0,-7.6) {\textbf{Pool / group structure}\\Sphere objects as logical devices; trioles/routed families as grouped device families};

\draw[arr] (visible) -- (kernel);
\draw[arr] (kernel) -- (os);
\draw[arr] (os) -- (store);
\draw[arr] (store) -- (pool);
\draw[arr] (visible.east) -- (public.west);
\draw[arr] (store.east) -- (internal.west);

\node[draw, rounded corners=5pt, fit=(visible)(kernel)(os)(store)(pool), inner sep=9pt, line width=0.6pt] (stackfit) {};
\node[above=0.15cm of stackfit.north, fill=white, inner sep=2pt] {\textbf{QAM as VM/OS/ZFS stack}};
\end{tikzpicture}%
}\caption{Graphical reading of the current QAM export as a virtual-machine/OS stack with a ZFS-style structural store. Visible tuple surfaces provide the public theorem interface; exposed payload kernels and kernel normal forms provide the later machine-state identity; snapshots, clones, search keys, and audits manage the internal proof store without changing the visible public QAM syntax.}
\label{fig:qam-vm-os-zfs}
\end{figure}

\begin{remark}[Why the ZFS analogy is structurally useful here]\label{rem:qam-zfs-utility}
The ZFS-style reading is not meant as decorative metaphor. It isolates four practical gains already visible in the present file. First, equal kernel classes can later be treated as mirror-like state reuse rather than as repeatedly independent visible tuple surfaces. Second, the deferred kernel-first path already behaves like copy-on-write: experimental compression statements are developed beside the active tuple proofs without rewriting the active line. Third, checkpoints, audits, and activation criteria behave like snapshots and scrub rules for the theorem store, since they record exactly which deferred branches are stable and which still remain non-authoritative. Fourth, once kernel normal forms and search keys are in place, repeated branch/shell/routed return families can be indexed and deduplicated as one structural graph beneath a readable tree-shaped theorem presentation. This is precisely why the systems analogy deserves a more prominent place next to the public QAM extraction itself.
\end{remark}

\begin{remark}[Current scope of the VM/OS/ZFS view]\label{rem:qam-vm-os-scope}
In the inherited v3.16.0 public-shape layer this systems view remains explanatory and organizational. It does not introduce a second formal semantics, and it does not yet activate RAID-like reconstruction as an active theorem rule. The present mathematically justified gains at that inherited layer are the mirror-like reuse of equal kernel classes, the snapshot/clone logic of active versus deferred proof branches, and the content-addressed search-key discipline attached to kernel normal forms. The v3.34.0 extension below activates only a guarded store-and-array vocabulary over the already available QAM/MML/MMA interface. It still does not turn mirror, stripe, parity, or repair language into an OP-closure certificate.
\end{remark}


